\documentclass[10pt,a4paper]{amsart}
\usepackage{amsmath,amssymb,amsthm,graphicx}
\usepackage{verbatim}
\newenvironment{reference}{\comment}{\endcomment}
%\newenvironment{reference}{}{}
\newenvironment{slogan}{\comment}{\endcomment}
\newenvironment{history}{\comment}{\endcomment}

% For commutative diagrams we use Xy-pic
\usepackage[all]{xy}

% We use 2cell for 2-commutative diagrams.
\xyoption{2cell}
% Theorem environments.
%
\theoremstyle{plain}
\newtheorem{theorem}[subsection]{Theorem}
\newtheorem{proposition}[subsection]{Proposition}
\newtheorem{lemma}[subsection]{Lemma}

\theoremstyle{definition}
\newtheorem{definition}[subsection]{Definition}
\newtheorem{example}[subsection]{Example}
\newtheorem{exercise}[subsection]{Exercise}
\newtheorem{situation}[subsection]{Situation}

\theoremstyle{remark}
\newtheorem{remark}[subsection]{Remark}
\newtheorem{remarks}[subsection]{Remarks}

\numberwithin{equation}{subsection}

% Macros
%
\def\lim{\mathop{\mathrm{lim}}\nolimits}
\def\colim{\mathop{\mathrm{colim}}\nolimits}
\def\Spec{\mathop{\mathrm{Spec}}}
\def\Hom{\mathop{\mathrm{Hom}}\nolimits}
\def\Ext{\mathop{\mathrm{Ext}}\nolimits}
\def\SheafHom{\mathop{\mathcal{H}\!\mathit{om}}\nolimits}
\def\SheafExt{\mathop{\mathcal{E}\!\mathit{xt}}\nolimits}
\def\Sch{\mathit{Sch}}
\def\Mor{\mathop{\mathrm{Mor}}\nolimits}
\def\Ob{\mathop{\mathrm{Ob}}\nolimits}
\def\Sh{\mathop{\mathit{Sh}}\nolimits}
\def\NL{\mathop{N\!L}\nolimits}
\def\CH{\mathop{\mathrm{CH}}\nolimits}
\def\proetale{{pro\text{-}\acute{e}tale}}
\def\etale{{\acute{e}tale}}
\def\QCoh{\mathit{QCoh}}
\def\Ker{\mathop{\mathrm{Ker}}}
\def\Im{\mathop{\mathrm{Im}}}
\def\Coker{\mathop{\mathrm{Coker}}}
\def\Coim{\mathop{\mathrm{Coim}}}

% Boxtimes
%
\DeclareMathSymbol{\boxtimes}{\mathbin}{AMSa}{"02}

%
% Macros for moduli stacks/spaces
%
\def\QCohstack{\mathcal{QC}\!\mathit{oh}}
\def\Cohstack{\mathcal{C}\!\mathit{oh}}
\def\Spacesstack{\mathcal{S}\!\mathit{paces}}
\def\Quotfunctor{\mathrm{Quot}}
\def\Hilbfunctor{\mathrm{Hilb}}
\def\Curvesstack{\mathcal{C}\!\mathit{urves}}
\def\Polarizedstack{\mathcal{P}\!\mathit{olarized}}
\def\Complexesstack{\mathcal{C}\!\mathit{omplexes}}
% \Pic is the operator that assigns to X its picard group, usage \Pic(X)
% \Picardstack_{X/B} denotes the Picard stack of X over B
% \Picardfunctor_{X/B} denotes the Picard functor of X over B
\def\Pic{\mathop{\mathrm{Pic}}\nolimits}
\def\Picardstack{\mathcal{P}\!\mathit{ic}}
\def\Picardfunctor{\mathrm{Pic}}
\def\Deformationcategory{\mathcal{D}\!\mathit{ef}}
\usepackage{cleveref}
\crefname{originallemma}{original lemma}{original lemmas}
\Crefname{originallemma}{Original lemma}{Original lemmas}
\makeatletter
\newlabel{algebra-lemma-strictly-henselian-functorial-improve}{{08HV}{source}{}{sourcepointer-1}{}}
\newlabel{algebra-lemma-strictly-henselian-functorial-improve@cref}{{[originallemma][0][]08HV}{[1][0][]source}{}{}{}}
\makeatother
\crefname{originalequation}{original equation}{original equations}
\Crefname{originalequation}{Original equation}{Original equations}
\makeatletter
\newlabel{equation-bc}{{09LQ}{source}{}{sourcepointer-2}{}}
\newlabel{equation-bc@cref}{{[originalequation][0][]09LQ}{[1][0][]source}{}{}{}}
\makeatother
\crefname{originalreference}{original reference}{original references}
\Crefname{originalreference}{Original reference}{Original references}
\makeatletter
\newlabel{lemma-derived-quasi-coherent-small-etale-site}{{071Q}{source}{}{sourcepointer-3}{}}
\newlabel{lemma-derived-quasi-coherent-small-etale-site@cref}{{[originalreference][0][]071Q}{[1][0][]source}{}{}{}}
\makeatother
\crefname{originallemma}{original lemma}{original lemmas}
\Crefname{originallemma}{Original lemma}{Original lemmas}
\makeatletter
\newlabel{more-algebra-lemma-dumb-properties-henselization}{{07QM}{source}{}{sourcepointer-4}{}}
\newlabel{more-algebra-lemma-dumb-properties-henselization@cref}{{[originallemma][0][]07QM}{[1][0][]source}{}{}{}}
\makeatother
\crefname{originallemma}{original lemma}{original lemmas}
\Crefname{originallemma}{Original lemma}{Original lemmas}
\makeatletter
\newlabel{more-algebra-lemma-lemma-tor-independent-flat-compare}{{0DJD}{source}{}{sourcepointer-5}{}}
\newlabel{more-algebra-lemma-lemma-tor-independent-flat-compare@cref}{{[originallemma][0][]0DJD}{[1][0][]source}{}{}{}}
\makeatother
\crefname{originallemma}{original lemma}{original lemmas}
\Crefname{originallemma}{Original lemma}{Original lemmas}
\makeatletter
\newlabel{perfect-lemma-single-complex-base-change-condition}{{0DJ8}{source}{}{sourcepointer-6}{}}
\newlabel{perfect-lemma-single-complex-base-change-condition@cref}{{[originallemma][0][]0DJ8}{[1][0][]source}{}{}{}}
\makeatother
\crefname{originalreference}{original reference}{original references}
\Crefname{originalreference}{Original reference}{Original references}
\makeatletter
\newlabel{remark-match-total-direct-images}{{08GH}{source}{}{sourcepointer-7}{}}
\newlabel{remark-match-total-direct-images@cref}{{[originalreference][0][]08GH}{[1][0][]source}{}{}{}}
\makeatother
\crefname{originallemma}{original lemma}{original lemmas}
\Crefname{originallemma}{Original lemma}{Original lemmas}
\makeatletter
\newlabel{spaces-properties-lemma-describe-etale-local-ring}{{04KF}{source}{}{sourcepointer-8}{}}
\newlabel{spaces-properties-lemma-describe-etale-local-ring@cref}{{[originallemma][0][]04KF}{[1][0][]source}{}{}{}}
\makeatother
\crefname{originallemma}{original lemma}{original lemmas}
\Crefname{originallemma}{Original lemma}{Original lemmas}
\makeatletter
\newlabel{spaces-properties-lemma-stalk-quasi-coherent}{{05VP}{source}{}{sourcepointer-9}{}}
\newlabel{spaces-properties-lemma-stalk-quasi-coherent@cref}{{[originallemma][0][]05VP}{[1][0][]source}{}{}{}}
\makeatother
\title{The base change isomorphism condition for a single complex is etale local and descends to schemes}
\author{The Stacks Project Authors}
\date{stacks-project:spaces-perfect.tex:local-snapshot-20261012; GNU FDL 1.2 or later}
\begin{document}
\maketitle
\noindent\textit{Editorial note.} All delivered mathematical text is the author's own, copied from the pinned Stacks Project chapter file spaces-perfect.tex; no statement or proof step was rewritten, reordered or removed.\par\medskip
\noindent\textit{Editorial note.} The counted proof is the author's complete argument: an etale-local reduction to the scheme case, a stalk computation through strict henselizations, and a faithful flatness argument identifying the comparison arrow as a base change, all inside the delivered span\par\medskip
\noindent\textit{Editorial note.} The label lemma-single-complex-base-change-condition names the counted statement, and the proof environment immediately following it in the source is delivered verbatim as its proof.\par\medskip
\noindent\textit{Editorial note.} The counted statement is carried by the author's own lemma environment in the chapter source.\par\medskip
\noindent\textit{Allowed standard prerequisites (editorial).}\begin{itemize}
\item The scheme-level base change criterion of Derived Categories of Schemes cited in the statement
\item Derived pullback to the small etale site and the strict henselization description of stalks of quasi-coherent complexes on algebraic spaces
\item Faithful flatness of the map from a local ring to its strict henselization, and the Tor-comparison lemma for strict henselizations of fibre products
\end{itemize}\medskip
\ifdefined\hypertarget\hypertarget{sourcepointer-1}{}\fi
\noindent\textit{Source reference.} Original lemma 08HV: sources/stacks-chow/tags.txt, source lines 10793--10793. Prerequisite pointer invoked at proof line(s) 5476; it names a standard fact of a sibling chapter that the counted argument uses, and every step of the counted proof built on it is retained verbatim in the delivered body.\par\medskip
\ifdefined\hypertarget\hypertarget{sourcepointer-2}{}\fi
\noindent\textit{Source reference.} Original equation 09LQ: sources/stacks-spaces-perfect/pinned-original-source.tex, source lines 5351--5351. Prerequisite pointer invoked at proof line(s) 5379; it names a standard fact of the same chapter that the counted argument uses, and every step of the counted proof built on it is retained verbatim in the delivered body.\par\medskip
\ifdefined\hypertarget\hypertarget{sourcepointer-3}{}\fi
\noindent\textit{Source reference.} Original reference 071Q: sources/stacks-spaces-perfect/pinned-original-source.tex, source lines 402--402. Prerequisite pointer invoked at proof line(s) 5410; it names a standard fact of the same chapter that the counted argument uses, and every step of the counted proof built on it is retained verbatim in the delivered body.\par\medskip
\ifdefined\hypertarget\hypertarget{sourcepointer-4}{}\fi
\noindent\textit{Source reference.} Original lemma 07QM: sources/stacks-chow/tags.txt, source lines 9837--9837. Prerequisite pointer invoked at proof line(s) 5452; it names a standard fact of a sibling chapter that the counted argument uses, and every step of the counted proof built on it is retained verbatim in the delivered body.\par\medskip
\ifdefined\hypertarget\hypertarget{sourcepointer-5}{}\fi
\noindent\textit{Source reference.} Original lemma 0DJD: sources/stacks-chow/tags.txt, source lines 16973--16973. Prerequisite pointer invoked at proof line(s) 5473; it names a standard fact of a sibling chapter that the counted argument uses, and every step of the counted proof built on it is retained verbatim in the delivered body.\par\medskip
\ifdefined\hypertarget\hypertarget{sourcepointer-6}{}\fi
\noindent\textit{Source reference.} Original lemma 0DJ8: sources/stacks-chow/tags.txt, source lines 16968--16968. Prerequisite pointer invoked at proof line(s) 5392; it names a standard fact of a sibling chapter that the counted argument uses, and every step of the counted proof built on it is retained verbatim in the delivered body.\par\medskip
\ifdefined\hypertarget\hypertarget{sourcepointer-7}{}\fi
\noindent\textit{Source reference.} Original reference 08GH: sources/stacks-spaces-perfect/pinned-original-source.tex, source lines 896--896. Prerequisite pointer invoked at proof line(s) 5413; it names a standard fact of the same chapter that the counted argument uses, and every step of the counted proof built on it is retained verbatim in the delivered body.\par\medskip
\ifdefined\hypertarget\hypertarget{sourcepointer-8}{}\fi
\noindent\textit{Source reference.} Original lemma 04KF: sources/stacks-chow/tags.txt, source lines 5971--5971. Prerequisite pointer invoked at proof line(s) 5417; it names a standard fact of a sibling chapter that the counted argument uses, and every step of the counted proof built on it is retained verbatim in the delivered body.\par\medskip
\ifdefined\hypertarget\hypertarget{sourcepointer-9}{}\fi
\noindent\textit{Source reference.} Original lemma 05VP: sources/stacks-chow/tags.txt, source lines 7558--7558. Prerequisite pointer invoked at proof line(s) 5428; it names a standard fact of a sibling chapter that the counted argument uses, and every step of the counted proof built on it is retained verbatim in the delivered body.\par\medskip
\begin{lemma}
\label{lemma-single-complex-base-change-condition}
Let $S$ be a scheme. Let
$$
\xymatrix{
X' \ar[r]_{g'} \ar[d]_{f'} &
X \ar[d]^f \\
Y' \ar[r]^g &
Y
}
$$
be a cartesian diagram of algebraic spaces over $S$.
Let $K \in D_\QCoh(\mathcal{O}_X)$ and let $L(g')^*K \to K'$
be a map in $D_\QCoh(\mathcal{O}_{X'})$. The following are equivalent
\begin{enumerate}
\item for any $x' \in X'$ and $i \in \mathbf{Z}$ the map (\ref{equation-bc})
is an isomorphism,
\item for any commutative diagram
$$
\xymatrix{
& U \ar[d] \ar[rd]^a \\
V' \ar[r] \ar[rd]^c & V \ar[rd]^b & X \ar[d]^f \\
& Y' \ar[r]^g & Y
}
$$
with $a, b, c$ \'etale, $U, V, V'$ schemes, and with $U' = V' \times_V U$
the equivalent conditions of
Derived Categories of Schemes, Lemma
\ref{perfect-lemma-single-complex-base-change-condition}
hold for $(U \to X)^*K$ and $(U' \to X')^*K'$, and
\item there is some diagram as in (2) with $U' \to X'$ surjective.
\end{enumerate}
\end{lemma}

\begin{proof}
Observe that (1) is \'etale local on $X'$. Working through formal
implications of what is known, we see that it suffices to prove
condition (1) of this lemma is equivalent to condition
(1) of Derived Categories of Schemes, Lemma
\ref{perfect-lemma-single-complex-base-change-condition}
if $X, Y, Y', X'$ are representable by schemes
$X_0, Y_0, Y'_0, X'_0$. Denote $f_0, g_0, g'_0, f'_0$ the
morphisms between these schemes corresponding to $f, g, g', f'$.
We may assume $K = \epsilon^*K_0$ and $K' = \epsilon^*K'_0$
for some objects $K_0 \in D_\QCoh(\mathcal{O}_{X_0})$
and $K'_0 \in D_\QCoh(\mathcal{O}_{X'_0})$, see
Lemma \ref{lemma-derived-quasi-coherent-small-etale-site}.
Moreover, the map $Lg^*K \to K'$ is the pullback
of a map $L(g_0)^*K_0 \to K'_0$ with notation as in
Remark \ref{remark-match-total-direct-images}.
Recall that $\mathcal{O}_{X, \overline{x}}$ is the
strict henselization of $\mathcal{O}_{X, x}$
(Properties of Spaces, Lemma
\ref{spaces-properties-lemma-describe-etale-local-ring})
and that we have
$$
K_{\overline{x}} =
K_{0, x} \otimes_{\mathcal{O}_{X, x}}^\mathbf{L} \mathcal{O}_{X, \overline{x}}
\quad\text{and}\quad
K'_{\overline{x}'} =
K'_{0, x'} \otimes_{\mathcal{O}_{X', x'}}^\mathbf{L}
\mathcal{O}_{X', \overline{x}'}
$$
(akin to Properties of Spaces, Lemma
\ref{spaces-properties-lemma-stalk-quasi-coherent}).
Consider the commutative diagram
$$
\xymatrix{
H^i(K_{\overline{x}}
\otimes_{\mathcal{O}_{Y, \overline{y}}}^\mathbf{L}
\mathcal{O}_{Y', \overline{y}'})
\otimes_{(\mathcal{O}_{X, \overline{x}}
\otimes_{\mathcal{O}_{Y, \overline{y}}}
\mathcal{O}_{Y', \overline{y}'})}
\mathcal{O}_{X', \overline{x}'}
\ar[r] &
H^i(K'_{\overline{x}'}) \\
H^i(K_{0, x} \otimes_{\mathcal{O}_{Y, y}}^\mathbf{L} \mathcal{O}_{Y', y'})
\otimes_{(\mathcal{O}_{X, x} \otimes_{\mathcal{O}_{Y, y}} \mathcal{O}_{Y', y'})}
\mathcal{O}_{X', x'}
\ar[u] \ar[r] &
H^i(K'_{0, x'}) \ar[u]
}
$$
We have to show that the lower horizontal arrow is an isomorphism if and only
if the upper horizontal arrow is an isomorphism. Since
$\mathcal{O}_{X', x'} \to \mathcal{O}_{X', \overline{x}'}$
is faithfully flat (More on Algebra, Lemma
\ref{more-algebra-lemma-dumb-properties-henselization})
it suffices to show that the top arrow is the base
change of the bottom arrow by this map.
This follows immediately from the relationships between
stalks given above for the objects on the right.
For the objects on the left it suffices to show that
\begin{align*}
&
H^i\left(
(K_{0, x}
\otimes_{\mathcal{O}_{X, x}}^\mathbf{L} \mathcal{O}_{X, \overline{x}})
\otimes_{\mathcal{O}_{Y, \overline{y}}}^\mathbf{L}
\mathcal{O}_{Y', \overline{y}'}\right) \\
& =
H^i(K_{0, x} \otimes_{\mathcal{O}_{Y, y}}^\mathbf{L} \mathcal{O}_{Y', y'})
\otimes_{(\mathcal{O}_{X, x} \otimes_{\mathcal{O}_{Y, y}} \mathcal{O}_{Y', y'})}
(\mathcal{O}_{X, \overline{x}}
\otimes_{\mathcal{O}_{Y, \overline{y}}}
\mathcal{O}_{Y', \overline{y}'})
\end{align*}
This follows from More on Algebra, Lemma
\ref{more-algebra-lemma-lemma-tor-independent-flat-compare}.
The flatness assumptions of this lemma hold by what was said
above as well as Algebra, Lemma
\ref{algebra-lemma-strictly-henselian-functorial-improve}
implying that
$\mathcal{O}_{X, \overline{x}}$ is the strict henselization of
$\mathcal{O}_{X, x} \otimes_{\mathcal{O}_{Y, y}}
\mathcal{O}_{Y, \overline{y}}$
and that
$\mathcal{O}_{Y', \overline{y}'}$ is the strict henselization of
$\mathcal{O}_{Y', y'} \otimes_{\mathcal{O}_{Y, y}}
\mathcal{O}_{Y, \overline{y}}$.
\end{proof}

\end{document}
